\chapter{绪论}

\section{研究背景}

《义务教育语文课程标准（2022年版）》把“文化自信”列为语文核心素养之首，强调让学生在真实的语言运用情境中感受中华文化的生命力，并在积极的语文实践活动中积累、建构和表现文化理解\cite{jiaoyubu2022}。综合性学习作为连接课堂学习与真实生活的重要载体，理应承担起引导学生认识地方文化、参与文化传承的责任。然而，在具体教学中，综合性学习常常被处理成资料搜集、手抄报制作或一次性的展示活动，文化资源与语言实践之间缺乏深度联结，学生的探究也容易停留在浅层信息的搬运上。

泉州于2021年以“泉州：宋元中国的世界海洋商贸中心”列入《世界遗产名录》，二十二处遗产点共同呈现了宋元时期海洋商贸、多元社群、城市营建与宗教信仰的历史图景\cite{quanzhou2021}。这些遗产不仅是地方文化的见证，也是学生身边可感、可访、可说的课程资源。开元寺的东西塔、洛阳桥的筏形基础、蟳埔村的蚵壳厝、九日山的祈风石刻，都蕴含着丰富的历史信息、文学意象与审美经验。如何把这些资源转化为语文学习内容，而不是简单叠加为“文化知识介绍”，是地方课程开发需要回答的问题。

项目化学习强调以真实问题驱动、以持续探究为主线、以公开成果为指向，与综合性学习所倡导的自主、合作、探究具有内在一致性\cite{xiaxuemei2018}。将项目化学习引入地方世界遗产教育，有可能把“参观—介绍”式的活动改造成“问题—探究—表达—反思”的学习链条，使学生在完成真实任务的过程中发展语言运用、思维品质、审美创造与文化理解。基于这一判断，本研究尝试在初中语文综合性学习中构建“世遗泉州”项目化学习方案，并通过行动研究检验其可行性。

\section{研究意义}

\subsection{理论意义}

已有研究对语文综合性学习的课程价值、资源开发与实施策略进行了较多讨论，但将地方世界遗产作为系统性课程资源，并以项目化学习组织语文实践的研究仍不充分。本研究尝试把文化回应性教学、建构主义学习观与逆向教学设计结合起来，分析地方文化资源转化为语文学习任务的内在机制，为语文综合性学习的资源开发提供一个新的分析框架。同时，本研究关注项目化学习在语文学科中的适切性问题，讨论“项目”如何服务于语言文字学习，而非替代语言文字学习，这有助于丰富项目化学习的学科化研究。

\subsection{实践意义}

对教师而言，本研究提供了一套可操作的项目设计流程、任务支架与评价工具，能够为地方文化主题的综合性学习提供参考。对学生而言，围绕身边遗产展开的探究有助于缩短语文学习与生活经验的距离，使学生在走访、访谈、阅读、写作与展示中形成对家乡文化的理解与认同。对学校而言，研究成果可以为校本课程建设、研学实践与校园文化活动提供可迁移的经验。需要说明的是，本研究的调查数据与教学案例均为示例性材料，正式使用时须依据实际教学情境重新采集与分析。

\section{研究现状}

\subsection{世界遗产教育研究}

世界遗产教育兴起于二十世纪后期，联合国教科文组织持续倡导通过教育增进青少年对遗产价值的理解与保护意识\cite{unesco2012}。相关研究大体形成三条路径：一是以遗产保护知识为主的课程渗透，二是以研学旅行为载体的体验学习，三是以地方认同为目标的社区参与。国内研究多关注遗产教育的价值阐释、课程资源开发与研学线路设计，研究对象以高校、博物馆和中小学综合实践活动为主\cite{shanjixiang2009}。总体来看，已有研究重视遗产知识的传递与保护意识的培养，但对遗产资源如何进入语文学科、如何转化为语言实践活动、如何评价学生在其中的语言发展，讨论仍较有限。

在地方层面，泉州世界遗产的研究多集中于历史学、建筑学与旅游管理领域，关注遗产价值、保护管理与旅游开发。教育领域的研究则常以“海丝文化进校园”“闽南文化校本课程”为主题，偏重活动介绍与经验总结，较少采用实证方法考察学生的学习过程与学习结果。这为本研究留下了空间：既需要尊重遗产研究的事实基础，也需要从语文学科的立场重新筛选、组织和转化这些资源。

\subsection{语文综合性学习研究}

自2001年课程改革设置综合性学习以来，相关研究大致经历了从理念倡导到实践探索、再到素养导向深化的过程。研究者普遍认为，综合性学习具有综合性、实践性、开放性等特征，应突破单篇课文与单一课时的局限，引导学生在问题解决中综合运用听说读写能力\cite{wangrongsheng2003}。2022年版课程标准进一步提出“以学习任务为载体，整合学习内容、情境、方法和资源等要素”，为综合性学习的设计提供了更明确的课程依据\cite{jiaoyubu2022}。

从实施现状看，综合性学习仍面临一些突出问题：一是目标泛化，活动热闹但语文学习目标不清晰；二是资源碎片化，地方文化材料被当作背景知识补充，未形成结构化的学习内容；三是过程指导不足，教师布置任务多、提供支架少；四是评价虚化，重成果展示而轻过程证据。已有研究提出的改进策略主要包括明确核心任务、创设真实情境、提供学习支架、采用表现性评价等，但这些策略如何在一线课堂中落地，仍需要更多基于设计的研究。

\subsection{项目化学习研究}

项目化学习源于杜威“做中学”的教育思想，经克伯屈“项目方法”的系统阐述后逐渐成为重要的教学设计取向\cite{kilpatrick1918}。巴克教育研究所将项目化学习概括为围绕复杂、真实的问题展开，学生通过持续探究、合作交流与成果创作获得知识与能力的过程\cite{buck2008}。近年来，国内学者强调项目化学习不应停留在“做产品”，而应指向学科核心概念与关键能力的理解\cite{xiaxuemei2020}。在语文领域，项目化学习被用于整本书阅读、写作教学、综合性学习与跨学科主题学习，形成了不少实践案例。

现有研究肯定了项目化学习在激发学习动机、促进合作探究、连接真实生活方面的价值，同时也指出其风险：项目主题过大容易挤压语言文字学习，任务设计不当容易造成“热闹的低效”，评价标准模糊则难以判断学生是否真正学会了语文。因此，项目化学习进入语文学科，需要回答“项目如何承载语文”“成果如何体现语文”“评价如何指向语文”三个关键问题。

\subsection{研究述评}

综合来看，世界遗产教育、语文综合性学习与项目化学习三个领域分别积累了较为丰富的研究成果，但三者之间的交叉研究仍显不足。具体表现为：第一，地方世界遗产资源进入语文课程的方式较为零散，缺乏从课程目标到学习任务、再到评价证据的系统设计；第二，项目化学习在语文综合性学习中的应用多停留在案例描述，缺少对实施过程、学生表现与改进策略的持续追踪；第三，现有评价多关注最终成果，对学生文化理解、语言运用与合作探究的过程性证据关注不够。基于此，本研究尝试以泉州世界遗产为资源，以项目化学习为组织方式，开展两轮行动研究，重点考察方案设计、实施效果与改进路径。

\section{研究问题与方法}

\subsection{研究问题}

本研究围绕以下三个问题展开：

\begin{enumerate}[label=（\arabic*）,leftmargin=3em]
  \item 如何将泉州世界遗产资源转化为初中语文综合性学习的项目化学习方案？
  \item 该方案在实施过程中对学生的文化理解、语言运用、合作探究与问题解决产生了怎样的影响？
  \item 方案实施中存在哪些问题，如何通过行动研究的迭代加以改进？
\end{enumerate}

\subsection{研究方法}

本研究采用行动研究法，辅以文献研究、问卷调查、课堂观察、作品分析与半结构化访谈。行动研究强调在研究者的教学实践中发现问题、设计方案、实施改进并反思结果，适合考察一项教学方案在真实情境中的形成过程。研究在泉州市某中学八年级一个班开展，班级共有学生四十二人，教学时间共八周，分两轮实施。第一轮聚焦方案可行性与基本流程，第二轮针对第一轮暴露的问题调整任务支架与评价方式。

数据收集包括：前测与后测问卷各一次，问卷围绕文化理解、语言运用、合作探究与问题解决四个维度设计；课堂观察记录十六课时；学生过程性作品与最终成果共四十二套；学生访谈十二人次，教师访谈四人次。问卷数据采用统计软件进行分析，访谈资料通过编码归纳主题，作品依据评价量表进行评分。研究遵循自愿参与和匿名化原则，涉及学生姓名、学校信息的内容均作匿名处理。

\section{研究思路与创新点}

本研究按照“问题提出—理论建构—方案设计—行动实施—效果分析—反思改进”的思路展开。首先通过文献研究与前期调查明确研究问题与现实需求；其次在文化回应性教学、建构主义学习观与逆向教学设计的启发下，提炼设计原则；再次围绕“如何让更多人读懂泉州世遗”这一驱动性问题，设计项目任务、学习支架与评价方案；随后开展两轮行动研究，收集问卷、访谈、课堂观察与作品数据；最后分析实施效果，讨论存在问题并提出改进建议。

本研究的创新之处主要体现在两个方面。其一，尝试建立“遗产资源—语文任务—学习证据”的转化链条，使地方文化资源不再只是背景材料，而成为阅读、表达、探究与评价的内容载体。其二，将项目化学习的迭代机制引入语文综合性学习，通过两轮行动研究记录方案调整的过程，呈现“设计—实施—反思—再设计”的实践逻辑，为一线教师提供可参考的改进路径。
